\section{Conclusion}
\label{sec:conclusion}

This paper compared four empirical studies on the dual role of LLMs in computational text and media sciences and tested one of their claims in a small experiment. As simulated participants, instruction-tuned models can reproduce several classic behavioural findings, but they show a hyper-accuracy distortion and reflect the opinions of some groups much better than others, and prompting only partly corrects this (RQ1, RQ2). In the own replication with a 0.5B open model pair, instruction tuning alone did not cause the distortion, which is consistent with the view that it depends on model scale as well as on alignment. As writers, several current models produce convincing disinformation articles, although safety behaviour differs strongly between models (RQ3). As detectors, fine-tuned and prompted models can recognise machine-generated or propagandistic text only partly and under favourable conditions (RQ3, RQ4). Across studies, this paper identified a generation--detection asymmetry and a controllability gap, and it points to the combination, not tested in the reviewed studies, of LLM-generated disinformation and LLM-based propaganda detection as a concrete direction for future work.
